% Nouveau EGA V (cote 25): archive cote 25, site page 103 (PDF page 104).
% typescript. Draft compared: none.
% AI-transcribed by Claude Opus 5.5 (claude-opus-5-5, adaptive thinking, effort high), Message Batches API, batch ega25b (msgbatch_01A7n5YuXubUPNK1zsepi8RG), 2026-11-13; not yet checked by eye.
\scanpage{25/103}
\mspage{\hand{6}}
% NOTE: a single handwritten diagonal line runs across the whole page, from the top left of the text down to the bottom right edge; it seems to cancel the page. The text is transcribed below as live text.

et  un homomorphisme  $f^{\hand{*}}_{*}(\underline{L}) \longleftarrow \underline{E}'$  (tel que l'homomorphisme \unsure{associé}
$f^{*}(\underline{E}') \to \underline{L}$  soit surjectif), modulo  isomorphisme en $\underline{L}$.  On dira que
$g\colon P \to P'$  est \uline{de degré} n si $\underline{L} = g^{*}(\underline{O}_{P'}(1))$  est de degré n. Il suffit
évidemment de savoir déterminer les homomorphismes de degré n pour tout n
donné.  Noter que si g est de degré n,\inserted[typed]{et P est à fibres de dimension $\geq 1$} on aura nécessairement $n \geq 0$
(car  sur un corps k, \xout{$\underline{O}_P(n)$} si dim P $\geq 1$, $\underline{O}_P(n)$ n'est engendré par ses
sections que si $n \geq 0$); bien entendu, on peut \xout{aussi} se restreindre au cas
où P a ses fibres de dim $\geq 1$ (en procédant comme dans 1.2.).  Comme on a
\[ f_{*}(f^{*}(\underline{M})(n)) = \underline{M} \otimes f_{*}(\underline{O}_P(n)) = \underline{M} \otimes \underline{\mathrm{Sym}}^{n}(\underline{E}) \hand{,} \]
on voit que la détermination des \over{g}{S}-morphismes  $g\colon P \to P'$  est ramenée
à la détermination des couples $(\underline{M},u)$ à isom près, où $\underline{M}$  est un Module
inversible sur S et  $u\colon \underline{E}' \to \underline{M}\otimes\underline{\mathrm{Sym}}^{n}(\underline{E})$  un homomorphisme  tel que
l'homomorphisme correspondant  $f^{*}(\underline{E}') \to f^{*}(\underline{M})(n)$  soit un épimorphisme.
Donc  g détermine un premier invariant de nature globale sur S, savoir
$\mathrm{cl}(\underline{M}) \in \mathrm{Pic}(S)$, et cet invariant fixé par un $\underline{M}$ choisi,  les g correspon-%
dants correspondent à un certain sous-ensemble de l'ensemble quotient
$\mathrm{Hom}(\underline{E}',\underline{M}\otimes\underline{\mathrm{Sym}}^{n}(\underline{E}))/\mathrm{I}\,(S,\underline{O}_S)^{*}$,  le passage au quotient par le groupe
$\mathrm{I}\,(S,\underline{O}_S)^{*}$ correspondant  au\xout{x}  ``modulo isomorphisme'' dans la description
des $g\colon P \to P'$ via Modules inversibles  (NB  les endomorphismes resp. automor
phismes d'un $\underline{L}$ inversible sur un fibré projectif $\underline{P}$, correspondant aux
sections resp. sections\hand{$|$}inversibles de $\underline{O}_P$ sur P,  ou encore des sections
resp sections inversibles de $\underline{O}_S \hand{$\xrightarrow{\sim}$} f_{\hand{*}}(\underline{O}_P)$ sur S).

Cas particuliers: \hand{1)} n=0 , alors  il faut prendre  les  homomorphismes
$\underline{E}' \to \underline{M}$  qui sont \deleted{partout non nuls}\inserted{\unsure{surjectifs} i.e. partout non nuls}, modulo isomorphisme en $\underline{M}$ : on trouve
exactement les  morphismes  $g\colon P \to P'$  de la forme  hf , où  h est une
section de  P' sur S  (savoir celle déterminée par  le quotient inversible
$\underline{M}$ de $\underline{E}'$). Donc les \inserted{S-}\uline{morphismes de  degré} 0  \uline{de} P \uline{dans} P' \uline{sont les morphis-}%

\margin{103}
